\documentclass[10pt]{article}
\usepackage[margin=0.9in]{geometry}
\usepackage{amsmath,amssymb}
\usepackage[T1]{fontenc}
\usepackage[utf8]{inputenc}
\usepackage{lmodern}
\usepackage{microtype}
\usepackage{booktabs}
\usepackage{enumitem}
\usepackage[hidelinks]{hyperref}
\hypersetup{pdftitle={A Two-Frequency Sub-Millimetre Dish Search at 542 and 596 GHz: A Pre-Registered Kill Test for a Dark-Energy-Tied Gravity Framework}}
\setlist{nosep,leftmargin=1.3em}
\setlength{\parskip}{4pt}\setlength{\parindent}{0pt}
\sloppy
\begin{document}
\begin{center}
{\Large\bfseries A Two-Frequency Sub-Millimetre Dish Search at 542 and 596\,GHz:\\ A Pre-Registered Kill Test for a Dark-Energy-Tied Gravity Framework}\\[6pt]
{\small Working note, 2026-10-15, v1.0. AI-assisted research programme; not peer reviewed.}
\end{center}

\begin{abstract}\noindent
We propose a narrow-band dish-antenna search for a 0.5\,mm line at two fixed frequencies, 541.7\,GHz and 595.5\,GHz.

\textbf{The framework and its prediction.} We test a framework in which the galactic acceleration scale $a_0$ is tied to the dark-energy density, $\rho_{\rm DE}=4a_0^2/(Gc^2)$. In this framework, the extra gravitating component (``cold energy'') has no non-gravitational channel. The only quantum mass the framework's vacuum induces is $m=(\rho_{\rm DE}\hbar^3/c^3)^{1/4}$. That gives 2.240\,meV for the canonical value of $a_0$ and 2.463\,meV for the alternative value.

\textbf{What the search decides.} If the cold component did couple to light as a kinetically mixed vector field of this mass, a metal plate would radiate a line at exactly these frequencies.
\begin{itemize}
\item \textbf{A detection would falsify the framework.} The framework predicts no such line, and this prediction was pre-registered.
\item \textbf{A null result would close} the only portal through which the cold component could be converted to light at useful efficiency.
\end{itemize}

\textbf{Reach.} A 1\,m$^2$ plate with a warm receiver reaches a mixing $\varepsilon\approx7\times10^{-12}$ in one hour. A 10\,m$^2$ plate with an SIS receiver reaches $1.4\times10^{-13}$ in 30 days. No published search covers 1--3\,meV; the nearest reached $\chi\lesssim10^{-8}$ at 0.67--0.92\,meV.
\end{abstract}

\section{Why these two frequencies}
\begin{itemize}
\item \textbf{The law.} The framework's galactic law is $g_N=g^2/(g+a_0)$, with $a_0=\kappa c\sqrt{G\rho_{\rm DE}}$. The value $\kappa=\tfrac12$ is fitted, not derived. With $\kappa=\tfrac12$, $\rho_{\rm DE}c^2=4a_0^2/G$.
\item \textbf{The mass.} The only dimensionally available quantum mass is $m=(\rho_{\rm DE}\hbar^3/c^3)^{1/4}$. We scanned every integer-exponent combination of $(a_0,G,c,\hbar,\rho_{\rm DE})$, and no other combination gives a visible-band mass [1].
\item \textbf{Two footings.} The two values of $a_0$ quoted for this programme ($9.3603\times10^{-11}$ and $1.1312\times10^{-10}$\,m\,s$^{-2}$) give
$$m=2.240\ \text{meV}\ (\nu=541.7\ \text{GHz},\ \lambda=0.553\ \text{mm}),\qquad m=2.463\ \text{meV}\ (\nu=595.5\ \text{GHz},\ \lambda=0.503\ \text{mm}).$$
\item \textbf{Line width.} A field with the galactic velocity dispersion $\sigma\simeq230$\,km\,s$^{-1}$ gives a line width $\Delta\nu=\nu\sigma^2/c^2\simeq0.3$\,MHz.
\item \textbf{The search space.} The search is two narrow windows, not a mass scan.
\end{itemize}

\section{Signal and reach}
\begin{itemize}
\item \textbf{Signal power.} A flat conducting plate of area $A$ in a kinetically mixed vector field of density $\rho$ emits power [2]
$$P=\tfrac23\,\varepsilon^2\rho c^3 A,$$
normal to its surface. Here $\tfrac23$ averages over random polarization.
\item \textbf{Density used.} We take the framework's local cold density, $\rho=0.01\,M_\odot\,{\rm pc}^{-3}$ ($\rho c^2=6.1\times10^{-5}$\,J\,m$^{-3}$). This assumes the hypothesis under test: that the cold component \emph{is} the converting field.
\item \textbf{Example rate.} At $\varepsilon=10^{-12}$ and $A=10$\,m$^2$, $P=1.2\times10^{-19}$\,W, or about $3\times10^{7}$ photons per day.
\item \textbf{Noise.} The radiometer equation gives $\sigma_P=k_BT_{\rm sys}\sqrt{\Delta\nu/t}$.
\end{itemize}

\begin{center}
\begin{tabular}{llccc}
\toprule
plate & receiver & 1 h & 1 day & 30 days\\
\midrule
1\,m$^2$ & warm, 1000\,K & $7\times10^{-12}$ & $3\times10^{-12}$ & $1.4\times10^{-12}$\\
1\,m$^2$ & SIS, 100\,K & $2\times10^{-12}$ & $1.0\times10^{-12}$ & $4.5\times10^{-13}$\\
10\,m$^2$ & SIS, 100\,K & $7\times10^{-13}$ & $3\times10^{-13}$ & $1.4\times10^{-13}$\\
\bottomrule
\end{tabular}\\[2pt]
{\footnotesize $5\sigma$ reach in $\varepsilon$ (lane CFG797 [1]). Receiver temperatures are provisional.}
\end{center}

\section{Apparatus}
\begin{itemize}
\item \textbf{Plate and mirror.} Use a flat or spherical polished aluminium plate, flat to better than about $\lambda/20\approx25\,\mu$m, and a parabolic mirror focusing onto the feed horn. Existing millimetre-wave dish experiments use this geometry [3].
\item \textbf{Receiver.}
  \begin{itemize}
  \item An SIS heterodyne mixer from radio-astronomy band-8/9 technology covers both lines with one LO retune.
  \item The backend is a spectrometer with $\lesssim30$\,kHz channels.
  \item A room-temperature Schottky receiver suffices for the converter-window test.
  \end{itemize}
\item \textbf{Path.} The 557\,GHz water line lies between the two targets. Purge the optical path with dry nitrogen, or evacuate it.
\item \textbf{Shielding and checks.}
  \begin{itemize}
  \item Use an absorber-lined enclosure and a reflector-in/out chop.
  \item Inject a calibration tone at each target frequency.
  \item A real dark-field line does not move when the plate is rotated about its normal. It does change with the plate's area and orientation.
  \end{itemize}
\item \textbf{Blinding.} Fix both frequency windows in advance (this note). Offset the channel labels until the analysis is frozen.
\end{itemize}

\section{What each outcome means}
\begin{itemize}
\item \textbf{A line at 541.7 or 595.5\,GHz, with the expected width and plate scaling.}
  \begin{itemize}
  \item The framework's statement that cold energy has no visible channel is falsified [1]. The framework as committed is dead.
  \item A dark sector with a light-coupled vector would be discovered.
  \end{itemize}
\item \textbf{No line, at $\varepsilon\lesssim10^{-12}$.}
  \begin{itemize}
  \item The framework's null prediction survives.
  \item The only portal that could convert the cold component to light at useful efficiency is closed. That window needs $\varepsilon\approx4$--$6\times10^{-9}$ [1]. Even a one-hour warm-receiver run covers it by more than $500\times$.
  \end{itemize}
\item \textbf{A line at any other frequency.} It is unrelated to this framework's prediction, but interesting in its own right.
\end{itemize}

\section{Caveats}
\begin{itemize}
\item The target mass relies on the identification $\rho_{\rm DE}=4a_0^2/(Gc^2)$, which is equivalent to the fitted coefficient $\kappa=\tfrac12$.
\item The local density and receiver temperatures are provisional.
\item Cosmological and stellar bounds on a 2.2--2.5\,meV kinetically mixed vector were checked only against the solar relation ($\varepsilon\lesssim6\times10^{-9}$) [1]. CMB-era bounds were not re-audited.
\item The comparison with published searches comes from a literature search and is provisional.
\end{itemize}

\vspace{4pt}\hrule\vspace{3pt}
{\footnotesize
[1] This programme (public repository): ARGON03 haloscope audit (native mass, kill test); lanes CFG723 (no visible channel), CFG795 (converter design space) and CFG797 (this forecast), each with frozen criteria and scripts.
[2] D.~Horns, J.~Jaeckel, A.~Lindner, A.~Lobanov, J.~Redondo, A.~Ringwald, JCAP 04 (2013) 016.
[3] S.~Knirck et al., ``First results from a hidden photon dark matter search in the meV sector using a plane-parabolic mirror system'', JCAP 11 (2018) 031.
AI assistance: Anthropic's Claude; parallel AI sessions.}
\end{document}
